\documentclass[11pt]{article}

\usepackage[utf8]{inputenc}
\usepackage[T1]{fontenc}
\usepackage{amsmath,amssymb,amsfonts}
\usepackage{graphicx}
\usepackage{hyperref}
\usepackage{geometry}
\geometry{margin=1in}

\title{Proyecto Tesis}
\author{Tomas Sierra Sanchez \\
\small Advisor: Valérie Gauthier Umaña \\
\small Co-Advisor: Juan Fernando Perez \\
\small Universidad de Los Andes, Bogotá, Colombia \\
\small Bachelor's in Computing and Systems Engineering}
\date{2026}

\begin{document}

\maketitle

% =====================================================
\section{State of the Art}
% =====================================================

\subsection{Differential Privacy}
% - RSA/ECC security based on integer factoring / discrete log
% - Shor's algorithm threat, NIST push for PQC standards
% - SDP definition: find low Hamming weight x such that Hx^T = y^T
% - SDitH v1.1 standard-form parity check matrix H = (H' | I_{n-k})
% - Reconstruction from partial vector x = (x_A | x_B)

\subsection{Quantum Computing}
% - Origin: Ishai et al., converting MPC protocols into ZKPs
% - Prover simulates protocol among N parties
% - SDitH v1.1 variants:
%   - Hypercube variant: D-dimensional geometric structure
%   - Threshold variant: Shamir-secret-based comparison scheme,
%     polynomial consistency check S(X)*Q(X) = P(X)*F(X)

\subsection{Quantum Machine Learning}
% - Cold-boot attacks and DRAM data remanence
% - Bit decay to a "base" state after power-off
% - Modelled as an asymmetric binary channel (noisy recovery)

% =====================================================
\section{Description}
% =====================================================
% - Focus: security of SDitH under BBLM, seed recovery for deterministic key expansion
% - Approach: understand scheme -> understand attack -> adapt attack to SDitH structure
% - Deliverable: algorithmic model based on BBLM to recover the secret seed
% - Target audience: specialists (concepts referenced, not re-derived)
% - Contribution to NIST PQC signature standardization evaluation of SDitH

% =====================================================
\section{Objectives}
% =====================================================

\subsection{General Objective}



\subsection{Specific Objectives}
\begin{itemize}
    \item % Understand SDitH (SDP-based signature scheme) and BBLM (seed-recovery attack)
    \item % Investigate vulnerabilities in SDitH's signature/key-generation from the secret seed
    \item % Build an attack algorithm to recover the seed generating SDitH's shared vectors/partitions
    \item % Determine whether SDitH v1.1 is a viable post-quantum signature scheme
\end{itemize}

% =====================================================
\section{Expected Results}
% =====================================================
% - Article on SDitH's security assessment against BBLM
% - Possible extension of the recovery algorithm to other MPCitH schemes
% - Simulation of both the signature scheme and the attack, using/adapting
%   published algorithms to SDitH's architecture
% - Two possible outcomes:
%   1. Attack succeeds -> demonstrate exploitable weakness for future patching
%   2. Attack fails -> support SDitH's viability as a PQC signature scheme

% =====================================================
\section{Ethical Considerations}
% =====================================================
% - Governed by professional ethics / defensive cybersecurity principles
% - Perspective of implementer/security evaluator, not attacker in the wild
% - Goals:
\begin{enumerate}
    \item % Strengthen the scheme: provide data for countermeasures/redundancy mitigation
    \item % Inform standardization: support NIST's PQC selection process
    \item % Prevent risks: controlled disclosure vs. post-deployment exploitation
\end{enumerate}
% - Argument: no harm to end users; standard, recommended cryptanalytic practice
% - Framed as "Adversary Modeling" to keep trust in new PQC standards conservative
%   and mathematically grounded

% =====================================================
\section{Activity Calendar}
% =====================================================
% 
% =====================================================
\section*{References}
% =====================================================
\begin{thebibliography}{}
\end{thebibliography}

\end{document}
